\NSPage{055}%
For a tuple \NSi{NS-F-P055-001}{V} of physical fields, we write
\NSFormula{NS-F-P055-002}{display}{5.23}
\begin{equation}
 |V|_m=\max_{|\alpha|+b\leq m}|\partial_x^\alpha\partial_t^bV|.
 \tag{5.23}\label{eq:5.23}
\end{equation}
This is the pointwise maximum of its Cartesian space--time derivatives through order
\NSi{NS-F-P055-003}{m}. We write
\NSi{NS-F-P055-004}{u_n=u_{r,n}e_r+u_{\theta,n}e_\theta+u_{z,n}e_z} for the physical velocity
coefficient in \eqref{eq:5.1}, and \NSi{NS-F-P055-005}{\mathcal T_0} for the leading stress profile of
\eqref{eq:4.11}. We define the finite tuple, before applying any cutoff in
\NSi{NS-F-P055-006}{q}, by
\NSFormula{NS-F-P055-007}{display}{none}
\[
 \mathcal U^{[N]}=
 (u_{\mathrm{slow}}^{[N]},p_{\mathrm{slow}}^{[N]},T_{\mathrm{slow}}^{[N]})
 =\sum_{n=0}^{N}(u_n,p_n,q^{-A-1/2+\lambda_n}\mathcal T_n).
\]
The profiles in this sum already include their radial cutoffs and moment corrections. For a
physical stress pair \NSi{NS-F-P055-008}{T=(T_\theta,T_z)}, we define
\NSFormula{NS-F-P055-009}{display}{5.24}
\begin{equation}
 \mathcal F_{\mathrm{slow}}(u,p,T)=\mathscr R(u,p)
 +(\partial_r+2/r)T_\theta e_\theta+(\partial_r+1/r)T_ze_z.
 \tag{5.24}\label{eq:5.24}
\end{equation}
The added terms are the tangential divergence of the stress retained for the later wave
construction.

\begin{proposition}\label{prop:5.3}
Fix the leading profiles of Theorem~\ref{thm:4.6}. There is a sequence of positive-order profiles
constructed at each order by Lemma~\ref{lem:5.1} and then Lemma~\ref{lem:5.2}, using the finalized
lower orders. Every finite velocity sum \NSi{NS-F-P055-010}{u_{\mathrm{slow}}^{[N]}} is divergence-free.
On every fixed compact profile range \NSi{NS-F-P055-011}{0\leq X\leq X_{\max}},
\NSi{NS-F-P055-012}{-1\leq\eta\leq1}, with \NSi{NS-F-P055-013}{0<X_{\max}<\infty}, for each
\NSi{NS-F-P055-014}{N,m\geq0} and \NSi{NS-F-P055-015}{0<q\leq1}, its augmented momentum residual
satisfies
\NSFormula{NS-F-P055-016}{display}{5.25}
\begin{equation}
 |\mathcal F_{\mathrm{slow}}(\mathcal U^{[N]})|_m
 \leq C_{N,m}q^{2h(N+1)-K_m},
 \qquad K_m\ \text{independent of }N.
 \tag{5.25}\label{eq:5.25}
\end{equation}
The constants may depend on the fixed profile range, and
\NSi{NS-F-P055-017}{C_{N,m}} may depend on \NSi{NS-F-P055-018}{N}; the exponent loss
\NSi{NS-F-P055-019}{K_m} does not.
\end{proposition}

\begin{proof}
\textbf{Step 1: construct the sequence and cancel the retained coefficients.}
At order \NSi{NS-F-P055-020}{n}, we apply Lemma~\ref{lem:5.1} with the already finalized orders
\NSi{NS-F-P055-021}{0,\ldots,n-1}. We then apply Lemma~\ref{lem:5.2} to extend and repair the
current profiles, reconstruct \NSi{NS-F-P055-022}{V_n,\Pi_n} from Equations~\eqref{eq:5.2}
and~\eqref{eq:5.5}, and define \NSi{NS-F-P055-023}{\mathcal T_n} by \eqref{eq:5.9}. These
operations preserve the equations on the retained inner interval and the regularity needed to
form \NSi{NS-F-P055-024}{\Omega_n/X} at the next order. Both lemmas use the same fixed radial
intervals, so the induction continues for all \NSi{NS-F-P055-025}{n}.

Equation~\eqref{eq:5.2} gives exact incompressibility at each order, hence for every finite sum.
The pressure equation holds globally. Differentiating the stress primitives gives
\NSFormula{NS-F-P055-026}{display}{none}
\[
 (\partial_R+2/R)\mathcal T_{n,\theta}=-\mathfrak r_{\theta,n},
 \qquad
 (\partial_R+1/R)\mathcal T_{n,z}=-\mathfrak r_{z,n}.
\]
Thus the tangential residual is exactly minus the stress divergence at each retained order.
The same leading identities hold at order zero by Proposition~\ref{prop:4.2}. After truncation at
\NSi{NS-F-P055-027}{N}, every uncancelled product or shifted viscous term has exponent increment
at least \NSi{NS-F-P055-028}{2h(N+1)}. For example, at \NSi{NS-F-P055-029}{N=1} the product of
two order-one velocity coefficients and axial viscosity applied to order one first enter at order
two.

\textbf{Step 2: estimate physical derivatives of the finite residual.}
The powers of \NSi{NS-F-P055-030}{q} introduced by physical differentiation follow directly from
the coordinates. For a profile \NSi{NS-F-P055-031}{g(x_\perp/\sqrt q,\eta)} smooth in its Cartesian
arguments on a compact profile range, a derivative with \NSi{NS-F-P055-032}{a=|\beta|} transverse,
\NSi{NS-F-P055-033}{k} axial and \NSi{NS-F-P055-034}{l} time differentiations satisfies
\NSFormula{NS-F-P055-035}{display}{5.26}
\begin{equation}
 \left|\partial_{x_\perp}^{\beta}\partial_z^k\partial_t^l
 \bigl[q^b g(x_\perp/\sqrt q,\eta)\bigr]\right|
 \leq C_m(1+|b|)^m\|g\|_{C^m}q^{b-a/2-Dk-l},
 \qquad m=a+k+l.
 \tag{5.26}\label{eq:5.26}
\end{equation}

\NSPage{056}%
This is the chain rule with
\NSi{NS-F-P056-001}{\partial_t=q^{-1}L^{-1}(-q\partial_q+X\partial_X+D\eta\partial_\eta)}
and
\NSi{NS-F-P056-002}{\partial_z=q^{-D}L^{-1}(2\eta q\partial_q-2\eta X\partial_X+d\partial_\eta)};
transverse differentiation introduces a factor \NSi{NS-F-P056-003}{q^{-1/2}} in these smooth
Cartesian profiles.

The powers of \NSi{NS-F-P056-004}{n} generated by differentiating
\NSi{NS-F-P056-005}{q^{\lambda_n}} enter the coefficient constants. The exponent reduction depends
only on the number and type of physical derivatives. On the common stress support,
\NSi{NS-F-P056-006}{X\geq X_a>0}, so
\NSi{NS-F-P056-007}{r^{-1}=q^{-1/2}(2X)^{-1/2}} also has derivative losses independent of
\NSi{NS-F-P056-008}{N}; near the axis the velocity is estimated through its smooth Cartesian
representative. There are only finitely many terms in a fixed truncated residual. Applying
\eqref{eq:5.26} to them gives \eqref{eq:5.25}, with
\NSi{NS-F-P056-009}{K_m} independent of \NSi{NS-F-P056-010}{N}.
\end{proof}

\subsection{Divergence-preserving summation of the coefficients.}\label{sec:5.4}
The finite tuples \NSi{NS-F-P056-011}{\mathcal U^{[N]}} have the improving residual decay in
\eqref{eq:5.25}, but the constants \NSi{NS-F-P056-012}{C_{n,m}} in their coefficient bounds
\eqref{eq:5.17} may grow too quickly for the full formal series \eqref{eq:5.1} to converge. We
obtain smooth fields by multiplying the \NSi{NS-F-P056-013}{n}\,th positive-order term by
\NSi{NS-F-P056-014}{\chi(c_nq)}, where \NSi{NS-F-P056-015}{\chi=1} near zero and
\NSi{NS-F-P056-016}{\chi=0} beyond one, and choosing \NSi{NS-F-P056-017}{c_n\to\infty}, as in
\eqref{eq:5.34}. Later terms then vanish outside successively smaller neighborhoods of
\NSi{NS-F-P056-018}{q=0}. On a compact set with \NSi{NS-F-P056-019}{q} bounded below, only
finitely many terms remain. We must show that the finite residual bound of Proposition~\ref{prop:5.3}
survives these cutoffs. This is the conclusion \eqref{eq:5.36} of Lemma~\ref{lem:5.4}, applied to the
augmented residual \eqref{eq:5.24}.

To preserve incompressibility, we apply these cutoffs to vector potentials before differentiating.
An axisymmetric Stokes streamfunction \NSi{NS-F-P056-020}{S} gives such a potential: if
\NSi{NS-F-P056-021}{S/r^2} is smooth in \NSi{NS-F-P056-022}{(r^2,z,t)}, then
\NSFormula{NS-F-P056-023}{display}{none}
\[
 \mathcal A=\frac Sr e_\theta=\frac S{r^2}(-x_2,x_1,0),
 \qquad
 \operatorname{curl}\mathcal A=-\frac{\partial_zS}{r}e_r+\frac{\partial_rS}{r}e_z.
\]
The assumed smoothness of \NSi{NS-F-P056-024}{S/r^2} as a function of
\NSi{NS-F-P056-025}{(r^2,z,t)} makes the potential smooth in Cartesian coordinates. For the
positive-order profiles constructed above, the physical streamfunctions and their vector potentials
are
\NSFormula{NS-F-P056-026}{display}{5.27}
\begin{equation}
 F_n=X\mathcal A_X(U_n),
 \qquad \mathcal S_n=q^{1-A+\lambda_n}F_n,
 \qquad \mathcal A_n=(\mathcal S_n/r)e_\theta,
 \qquad
 u_{z,n}=\partial_s\mathcal S_n,
 \qquad ru_{r,n}=-\partial_z\mathcal S_n.
 \tag{5.27}\label{eq:5.27}
\end{equation}
Here \NSi{NS-F-P056-027}{s=r^2/2}, as in \eqref{eq:4.1}. Indeed
\NSi{NS-F-P056-028}{\operatorname{curl}\mathcal A_n=u_{r,n}e_r+u_{z,n}e_z} gives the radial and
axial velocity components. In Cartesian coordinates, using \NSi{NS-F-P056-029}{r^2=2qX},
\NSFormula{NS-F-P056-030}{display}{none}
\[
 \mathcal A_n=\frac{\mathcal S_n}{r^2}(-x_2,x_1,0)
 =\frac12q^{-A+\lambda_n}\frac{F_n}{X}(-x_2,x_1,0).
\]
Since \NSi{NS-F-P056-031}{F_n/X} is smooth at \NSi{NS-F-P056-032}{X=0}, this expression extends
smoothly across the axis for \NSi{NS-F-P056-033}{q>0}. Lemma~\ref{lem:5.2} supplies common
supports and bounds on every fixed coefficient derivative, with no bound uniform in
\NSi{NS-F-P056-034}{n} required.

Cutting off \NSi{NS-F-P056-035}{\mathcal A_n} and taking its curl includes the extra velocity term
from differentiating the cutoff, so the result remains divergence-free. We formulate the summation
argument for tuples of physical fields, allowing its later application to the wave and mean correction
sequence. For the present application its indices are \NSi{NS-F-P056-036}{j=n}, its decay orders are
\NSi{NS-F-P056-037}{g_n=2nh}, and its entries are
\NSFormula{NS-F-P056-038}{display}{none}
\[
 A_n=\mathcal A_n,
 \qquad B_n=u_{\theta,n}e_\theta,
 \qquad p_n=q^{-2A+\lambda_n}\Pi_n,
 \qquad T_n=q^{-A-1/2+\lambda_n}\mathcal T_n.
\]
Its differential polynomial will be \NSi{NS-F-P056-039}{\mathcal F_{\mathrm{slow}}} from
\eqref{eq:5.24}. In the lemma, \NSi{NS-F-P056-040}{U_j} denotes a tuple of physical fields, not the
axial profile \NSi{NS-F-P056-041}{U_n(X,\eta)}. We use the derivative notation \eqref{eq:5.23} and
set \NSi{NS-F-P056-042}{\Lambda_{\log}(q)=1+|\log q|}. All estimates in the lemma hold uniformly
over its stated domain.

\NSPage{057}%
The general argument needs a common domain, increasing decay orders, and derivative losses
independent of the summation index. Let \NSi{NS-F-P057-001}{\Omega} be a domain on which
\NSi{NS-F-P057-002}{0<q<q_0\leq1}. Assume that \NSi{NS-F-P057-003}{q=q(z,t)} is smooth, is
bounded below on compact subsets of \NSi{NS-F-P057-004}{\Omega}, and satisfies
\NSFormula{NS-F-P057-005}{display}{5.28}
\begin{equation}
 |\partial_z^a\partial_t^bq|\leq C_{a,b}q^{1-aD-b},
 \qquad a,b\geq0,
 \qquad 0<D<1.
 \tag{5.28}\label{eq:5.28}
\end{equation}
Let \NSi{NS-F-P057-006}{U_0=(u_0,p_0,T_0)} be smooth, with
\NSi{NS-F-P057-007}{\operatorname{div}u_0=0}. Assume that for every \NSi{NS-F-P057-008}{m}
there are finite \NSi{NS-F-P057-009}{C_m,K_m,P_m} such that
\NSi{NS-F-P057-010}{|U_0|_m\leq C_mq^{-K_m}\Lambda_{\log}(q)^{P_m}}. The entries
\NSi{NS-F-P057-011}{T_0} and \NSi{NS-F-P057-012}{T_j} below may be omitted. For
\NSi{NS-F-P057-013}{j\geq1}, suppose that the smooth physical tuples
\NSFormula{NS-F-P057-014}{display}{5.29}
\begin{equation}
 Z_j=(A_j,B_j,p_j,T_j),
 \qquad B_j=b_j(r,z,t)e_\theta,
 \qquad U_j=(\operatorname{curl}A_j+B_j,p_j,T_j)
 \tag{5.29}\label{eq:5.29}
\end{equation}
are defined on this same domain. Require smooth Cartesian representatives at the axis and smooth
zero extensions at every lateral boundary of their spatial supports. Suppose that, for every
\NSi{NS-F-P057-015}{j,m}, there are finite constants
\NSi{NS-F-P057-016}{C_{j,m},P_{j,m}} such that
\NSFormula{NS-F-P057-017}{display}{5.30}
\begin{equation}
 |Z_j|_m\leq C_{j,m}q^{g_j-\ell_m}\Lambda_{\log}(q)^{P_{j,m}},
 \qquad 0<g_1\leq g_2\leq\cdots\longrightarrow\infty,
 \tag{5.30}\label{eq:5.30}
\end{equation}
where the nonnegative, nondecreasing losses \NSi{NS-F-P057-018}{\ell_m} are independent of
\NSi{NS-F-P057-019}{j}.

For each finite \NSi{NS-F-P057-020}{J}, we set
\NSi{NS-F-P057-021}{U^{[J]}=U_0+\sum_{j=1}^{J}U_j}. We write the scalar components of a tuple as
\NSi{NS-F-P057-022}{U=(U^1,\ldots,U^{N_U})}. A fixed differential polynomial of finite order and
degree has, componentwise, the form
\NSFormula{NS-F-P057-023}{display}{5.31}
\begin{equation}
 (\mathcal F(U))^a=c_{a0}(x,t)
 +\sum_{\nu=1}^{M_a}c_{a\nu}(x,t)
 \prod_{r=1}^{d_{a\nu}}\partial_{x,t}^{\beta_{a\nu r}}U^{k_{a\nu r}}.
 \tag{5.31}\label{eq:5.31}
\end{equation}
Here the output index runs through the fixed finite set \NSi{NS-F-P057-024}{1\leq a\leq N_F},
each \NSi{NS-F-P057-025}{M_a} is finite, and the fixed integers
\NSi{NS-F-P057-026}{s\geq0,d\geq1} bound the order and degree:
\NSi{NS-F-P057-027}{1\leq d_{a\nu}\leq d,\ |\beta_{a\nu r}|\leq s}, and
\NSi{NS-F-P057-028}{1\leq k_{a\nu r}\leq N_U}. For
\NSi{NS-F-P057-029}{\beta=(\alpha,b)\in\mathbb N_0^4},
\NSi{NS-F-P057-030}{\partial_{x,t}^{\beta}=\partial_x^\alpha\partial_t^b}. The coefficient
functions, indices, and order and degree bounds \NSi{NS-F-P057-031}{s,d} are fixed independently of
the summation and truncation indices.

For each nonconstant monomial in \eqref{eq:5.31}, fix a common region
\NSi{NS-F-P057-032}{\Omega_{a\nu}} containing the supports of its field factors in the base and in
every increment, both before and after applying cutoffs. For the term independent of
\NSi{NS-F-P057-033}{U}, we set \NSi{NS-F-P057-034}{\Omega_{a0}=\Omega}. For every coefficient
\NSi{NS-F-P057-035}{c_{a\nu}}, \NSi{NS-F-P057-036}{0\leq\nu\leq M_a}, and every Cartesian
space--time multi-index \NSi{NS-F-P057-037}{I\in\mathbb N_0^4}, require
\NSFormula{NS-F-P057-038}{display}{5.32}
\begin{equation}
 |\partial_{x,t}^{I}c_{a\nu}(x,t)|
 \leq C_Iq^{-K_I}\Lambda_{\log}(q)^{P_I}
 \qquad ((x,t)\in\Omega_{a\nu}).
 \tag{5.32}\label{eq:5.32}
\end{equation}
The regions are independent of truncation. Since the coefficient list is finite, the finite constants
\NSi{NS-F-P057-039}{C_I,K_I,P_I} may be chosen common to that list for each fixed
\NSi{NS-F-P057-040}{I}; they are independent of the truncation index. Thus the bounds hold
throughout a comparison of partial sums and cutoff sums. The terms
\NSi{NS-F-P057-041}{c_{a0}}, which are independent of \NSi{NS-F-P057-042}{U}, cancel from that
comparison. The following residual hypothesis concerns the full \NSi{NS-F-P057-043}{\mathcal F},
including those terms: assume that each finite partial sum satisfies
\NSFormula{NS-F-P057-044}{display}{5.33}
\begin{equation}
 |\mathcal F(U^{[J]})|_m
 \leq C_{J,m}q^{\rho_J-K_m^{\mathcal F}}\Lambda_{\log}(q)^{P_{J,m}}+E_{J,m},
 \qquad \rho_J\longrightarrow\infty,
 \tag{5.33}\label{eq:5.33}
\end{equation}
where \NSi{NS-F-P057-045}{K_m^{\mathcal F}} is independent of \NSi{NS-F-P057-046}{J} and, for every fixed
\NSi{NS-F-P057-047}{J,m,N},
\NSi{NS-F-P057-048}{0\leq E_{J,m}\leq C_{J,m,N}q^N}. These bounds, including the flat remainder,
are uniform over the full domain at that fixed stage.

\begin{lemma}\label{lem:5.4}
There are numbers \NSi{NS-F-P057-049}{a_{j+1}\geq2a_j}, with
\NSi{NS-F-P057-050}{a_1^{-1}<q_0}, and a fixed smooth cutoff \NSi{NS-F-P057-051}{\chi}, equal to one
on \NSi{NS-F-P057-052}{[0,1/2]} and zero on \NSi{NS-F-P057-053}{[1,\infty)}, such that
\NSFormula{NS-F-P057-054}{display}{5.34}
\begin{equation}
 U=U_0+\sum_{j\geq1}
 \bigl(\operatorname{curl}(\chi(a_jq)A_j)+\chi(a_jq)B_j,
       \chi(a_jq)p_j,\chi(a_jq)T_j\bigr)
 \tag{5.34}\label{eq:5.34}
\end{equation}

\NSPage{058}%
is locally finite and smooth, and its velocity component \NSi{NS-F-P058-001}{u} satisfies
\NSi{NS-F-P058-002}{\operatorname{div}u=0}. Its added terms extend smoothly by zero across the outer
boundary \NSi{NS-F-P058-003}{q=q_0}. For every \NSi{NS-F-P058-004}{m} and
\NSi{NS-F-P058-005}{J\geq\max(1,m)} one has
\NSFormula{NS-F-P058-006}{display}{5.35}
\begin{equation}
 |U-U^{[J]}|_m\leq2^{-J}q^{g_{J+1}/2-\ell_m'}
 \qquad\left(0<q<\frac1{2a_J}\right),
 \tag{5.35}\label{eq:5.35}
\end{equation}
where \NSi{NS-F-P058-007}{\ell_m'} is independent of \NSi{NS-F-P058-008}{J}.

Under the residual hypothesis \eqref{eq:5.33}, every Cartesian space--time derivative of
\NSi{NS-F-P058-009}{\mathcal F(U)} vanishes to infinite order as
\NSi{NS-F-P058-010}{q\downarrow0}:
\NSFormula{NS-F-P058-011}{display}{5.36}
\begin{equation}
 \forall m,N\qquad\exists\delta,C>0\qquad
 |\mathcal F(U)|_m\leq Cq^N
 \qquad(0<q<\delta).
 \tag{5.36}\label{eq:5.36}
\end{equation}
\end{lemma}

\begin{proof}
We bound the derivatives of each cutoff, choose the cutoffs so the tail is summable after any fixed
number of derivatives, and compare the residual with one finite partial sum before cutoffs. In these
bounds, each prescribed Cartesian derivative reduces the power of \NSi{NS-F-P058-012}{q} by an
amount independent of the truncation order.

\textbf{Step 1: bound the cutoff and curl derivatives.}
For a positive-order derivative of \NSi{NS-F-P058-013}{\chi(aq)}, the chain rule gives terms of the
form
\NSFormula{NS-F-P058-014}{display}{none}
\[
 a^k\chi^{(k)}(aq)\prod_{\nu=1}^{k}\partial_z^{a_\nu}\partial_t^{b_\nu}q,
 \qquad
 \sum_\nu a_\nu=a',
 \qquad
 \sum_\nu b_\nu=b'.
\]
On their support, \NSi{NS-F-P058-015}{1/2\leq aq\leq1}. Using \eqref{eq:5.28}, the factors
\NSi{NS-F-P058-016}{q^k} compensate for \NSi{NS-F-P058-017}{a^k}, since
\NSi{NS-F-P058-018}{aq} is bounded on the cutoff support. Consequently
\NSFormula{NS-F-P058-019}{display}{none}
\[
 |\partial_z^{a'}\partial_t^{b'}\chi(aq)|
 \leq C_{a',b'}q^{-a'D-b'},
\]
with constants independent of \NSi{NS-F-P058-020}{a}. Since
\NSi{NS-F-P058-021}{q=q(z,t)},
\NSi{NS-F-P058-022}{\partial_{x_1}q=\partial_{x_2}q=0}. In particular the commutator in
\NSFormula{NS-F-P058-023}{display}{none}
\[
 \operatorname{curl}(\chi(aq)A_j)
 =\chi(aq)\operatorname{curl}A_j+a\chi'(aq)\nabla q\times A_j
\]
is controlled by the assumed derivatives of the potentials. The product rule therefore bounds
derivatives through order \NSi{NS-F-P058-024}{m} of the \NSi{NS-F-P058-025}{j}\,th summand in
\eqref{eq:5.34} by
\NSi{NS-F-P058-026}{\widehat C_{j,m}q^{g_j-\ell_m'}\Lambda_{\log}(q)^{\widehat P_{j,m}}}, where
the exponent reduction \NSi{NS-F-P058-027}{\ell_m'} depends only on
\NSi{NS-F-P058-028}{m} and \NSi{NS-F-P058-029}{\ell_{m+1}}.

\textbf{Step 2: choose the cutoffs and estimate the tail.}
We choose the scales \NSi{NS-F-P058-030}{a_j} recursively so that
\NSi{NS-F-P058-031}{a_{j+1}\geq2a_j}, each cutoff \NSi{NS-F-P058-032}{\chi(a_jq)} is supported in
\NSi{NS-F-P058-033}{q<q_0}, and
\NSFormula{NS-F-P058-034}{display}{5.37}
\begin{equation}
 \widehat C_{j,m}\Lambda_{\log}(q)^{\widehat P_{j,m}}q^{g_j/2}
 \leq2^{-j}
 \qquad(0<q\leq a_j^{-1},\ 0\leq m\leq j).
 \tag{5.37}\label{eq:5.37}
\end{equation}
Only finitely many requirements are imposed at step \NSi{NS-F-P058-035}{j}, and each holds for all
sufficiently small \NSi{NS-F-P058-036}{q} because \NSi{NS-F-P058-037}{g_j>0}. Derivatives through
order \NSi{NS-F-P058-038}{m} of the resulting cutoff increment are bounded by
\NSi{NS-F-P058-039}{2^{-j}q^{g_j/2-\ell_m'}} whenever \NSi{NS-F-P058-040}{m\leq j}. On a compact
subset where \NSi{NS-F-P058-041}{q\geq\delta>0}, all terms with
\NSi{NS-F-P058-042}{a_j>\delta^{-1}} vanish. Thus the sum is locally finite. The curl of any vector
field is divergence-free, and \NSi{NS-F-P058-043}{\chi(a_jq)B_j} is divergence-free because its
angular amplitude is independent of \NSi{NS-F-P058-044}{\theta}. Smoothness and the stated support
extensions follow termwise.

For fixed \NSi{NS-F-P058-045}{m} and \NSi{NS-F-P058-046}{J\geq\max(1,m)}, when
\NSi{NS-F-P058-047}{q<(2a_J)^{-1}} the first \NSi{NS-F-P058-048}{J} cutoffs equal one. Since
\NSi{NS-F-P058-049}{q<1} and \NSi{NS-F-P058-050}{g_j} is nondecreasing, the remaining terms obey
\NSFormula{NS-F-P058-051}{display}{none}
\[
 \sum_{j>J}2^{-j}q^{g_j/2-\ell_m'}
 \leq2^{-J}q^{g_{J+1}/2-\ell_m'},
\]
which proves \eqref{eq:5.35}.

\NSPage{059}%
\textbf{Step 3: compare the residual with one finite partial sum.}
The cutoff sum \NSi{NS-F-P059-001}{U} is now smooth, and its velocity component is divergence-free.
To prove flatness of \NSi{NS-F-P059-002}{\mathcal F(U)}, we compare it with
\NSi{NS-F-P059-003}{\mathcal F(U^{[J]})} for a finite partial sum with sufficiently high residual
decay, and control the change caused by \NSi{NS-F-P059-004}{U-U^{[J]}}. The assumed bound on
\NSi{NS-F-P059-005}{|U_0|_m} and \eqref{eq:5.30} give
\NSFormula{NS-F-P059-006}{display}{5.38}
\begin{equation}
 |U^{[J]}|_m\leq C_{J,m}q^{-K_m}\Lambda_{\log}(q)^{P_{J,m}},
 \tag{5.38}\label{eq:5.38}
\end{equation}
with \NSi{NS-F-P059-007}{K_m\geq0} independent of \NSi{NS-F-P059-008}{J}: every increment has
\NSi{NS-F-P059-009}{g_j>0}, and taking the curl requires one additional derivative of its potential.
The order and degree bounds \NSi{NS-F-P059-010}{s,d} are those in \eqref{eq:5.31}. Subtracting two
monomials one factor at a time, then applying the product rule and \eqref{eq:5.32}, gives for
\NSi{NS-F-P059-011}{e=U-U^{[J]}}
\NSFormula{NS-F-P059-012}{display}{5.39}
\begin{equation}
 |\mathcal F(U)-\mathcal F(U^{[J]})|_m
 \leq C_mq^{-H_m}|e|_{m+s}
 \bigl(1+|U^{[J]}|_{m+s}+|e|_{m+s}\bigr)^{d-1},
 \tag{5.39}\label{eq:5.39}
\end{equation}
where \NSi{NS-F-P059-013}{H_m} is independent of \NSi{NS-F-P059-014}{J}; the finitely many
coefficient logarithms are absorbed into this fixed power of \NSi{NS-F-P059-015}{q^{-1}}. A
polynomial independent of \NSi{NS-F-P059-016}{U} has zero difference and needs no estimate.

For fixed derivative order \NSi{NS-F-P059-017}{m} and decay order \NSi{NS-F-P059-018}{N}, we
choose \NSi{NS-F-P059-019}{J\geq m+s} large enough that
\NSFormula{NS-F-P059-020}{display}{none}
\[
 \frac{g_{J+1}}2-\ell_{m+s}'
 \geq N+H_m+(d-1)(K_{m+s}+1),
 \qquad
 \rho_J-K_m^{\mathcal F}\geq N+1.
\]
We also take \NSi{NS-F-P059-021}{J} large enough that
\NSi{NS-F-P059-022}{g_{J+1}/2-\ell_{m+s}'} is nonnegative, and then keep
\NSi{NS-F-P059-023}{J} fixed. After shrinking the \NSi{NS-F-P059-024}{q}-neighborhood,
\NSi{NS-F-P059-025}{\chi(a_jq)=1} for \NSi{NS-F-P059-026}{j\leq J}, the constants and logarithms
in \eqref{eq:5.38} are absorbed by \NSi{NS-F-P059-027}{q^{-1}}, and
\NSi{NS-F-P059-028}{|e|_{m+s}\leq1} by \eqref{eq:5.35}. Formula~\eqref{eq:5.39} is then
\NSi{NS-F-P059-029}{O(q^N)}. The same holds for \eqref{eq:5.33}, after absorbing its logarithm by
one power and using the flatness of this fixed \NSi{NS-F-P059-030}{E_{J,m}}. This proves
\eqref{eq:5.36}. The order of choices is \NSi{NS-F-P059-031}{(m,N)}, then
\NSi{NS-F-P059-032}{J}, then the neighborhood and constants; the comparison uses only the
remainder \NSi{NS-F-P059-033}{E_{J,m}} at that fixed truncation.
\end{proof}

A finite block of initial terms with nonpositive decay orders may receive one common cutoff,
applied as in \eqref{eq:5.34}, and be included in the base \NSi{NS-F-P059-034}{U_0}. Its cutoff
equals one near \NSi{NS-F-P059-035}{q=0}, preserving every comparison with a finite partial sum
there.

For smooth vector fields \NSi{NS-F-P059-036}{v,e} and smooth scalar fields
\NSi{NS-F-P059-037}{p,\pi}, the physical residual comparison is
\NSFormula{NS-F-P059-038}{display}{none}
\[
 \mathscr R(v+e,p+\pi)-\mathscr R(v,p)
 =\partial_te-\Delta e+\nabla\pi+(v\mathbin\cdot\nabla)e+(e\mathbin\cdot\nabla)v
 +(e\mathbin\cdot\nabla)e.
\]
For \NSi{NS-F-P059-039}{\mathcal F_{\mathrm{slow}}}, we add the difference of the linear tangential
stress-divergence term in \eqref{eq:5.24}. On the common annular stress support, the cylindrical
coefficients and each of their prescribed Cartesian derivatives are bounded by a fixed negative power
of \NSi{NS-F-P059-040}{q}. In the final correction sequence,
\NSi{NS-F-P059-041}{\mathcal F=\mathscr R} and there is no stress variable.

The same cutoff choices can meet further normalized estimates when those estimates have been proved
for the coefficients. More precisely, at step \NSi{NS-F-P059-042}{j} one may add any finite list of
requirements
\NSFormula{NS-F-P059-043}{display}{5.40}
\begin{equation}
 B_{j,k}\Lambda_{\log}(q)^{P_{j,k}}q^{\gamma_{j,k}}\leq2^{-j},
 \qquad \gamma_{j,k}>0,
 \tag{5.40}\label{eq:5.40}
\end{equation}
to \eqref{eq:5.37}. For derivatives of the normalized expansion coefficients and boundary weights,
we impose these additional coefficient estimates using the corresponding cutoff product rules. To
recover a full expansion after a fixed truncation \NSi{NS-F-P059-044}{N} and at a fixed derivative
order \NSi{NS-F-P059-045}{m}, we choose
\NSi{NS-F-P059-046}{M\geq\max(N,m)} with
\NSi{NS-F-P059-047}{g_{M+1}/2\geq g_{N+1}}: the finite block from
\NSi{NS-F-P059-048}{N+1} through \NSi{NS-F-P059-049}{M} retains its original orders near
\NSi{NS-F-P059-050}{q=0}, while \eqref{eq:5.35} controls the remaining tail at order
\NSi{NS-F-P059-051}{g_{N+1}}, with its fixed exponent reduction
\NSi{NS-F-P059-052}{\ell_m'}.

\NSPage{060}%
\subsection{The realized base field.}\label{sec:5.5}
The finite tuples \NSi{NS-F-P060-001}{\mathcal U^{[N]}} from Proposition~\ref{prop:5.3} now satisfy
the residual estimates required by Lemma~\ref{lem:5.4}. In this subsection we apply that lemma to
the coefficient fields and their potentials \NSi{NS-F-P060-002}{\mathcal A_n} in \eqref{eq:5.27} to
construct the smooth background \NSi{NS-F-P060-003}{(u_B,p_B)}. We must also retain the leading
profile as in \eqref{eq:5.42}, the exact heat exterior \eqref{eq:4.29}, and the weighted stress bounds
\eqref{eq:5.43} used to choose the wave amplitudes and their corrections in Propositions~\ref{prop:7.5}
and~\ref{prop:7.6}.

\begin{proposition}\label{prop:5.5}
There are smooth axisymmetric fields \NSi{NS-F-P060-004}{(u_B,p_B)} for
\NSi{NS-F-P060-005}{q>0} and two physical tangential stress components
\NSi{NS-F-P060-006}{\mathcal T_{\mathrm{phys}}}, supported where
\NSi{NS-F-P060-007}{X_a\leq X(r,z,t)\leq X_b}, such that
\NSi{NS-F-P060-008}{\operatorname{div}u_B=0} and
\NSFormula{NS-F-P060-009}{display}{5.41}
\begin{equation}
 \mathscr R(u_B,p_B)
 =-(\partial_r+2/r)\mathcal T_{\mathrm{phys},\theta}e_\theta
  -(\partial_r+1/r)\mathcal T_{\mathrm{phys},z}e_z+\mathcal E_B.
 \tag{5.41}\label{eq:5.41}
\end{equation}
For every fixed \NSi{NS-F-P060-010}{0<X_{\max}<\infty}, uniformly on the profile rectangle
\NSi{NS-F-P060-011}{[0,X_{\max}]\times[-1,1]} in \NSi{NS-F-P060-012}{(X,\eta)}, for every
Cartesian space--time derivative \NSi{NS-F-P060-013}{\partial^\alpha} and every
\NSi{NS-F-P060-014}{M>0},
\NSi{NS-F-P060-015}{|\partial^\alpha\mathcal E_B|\leq C_{\alpha,M}q^M} as
\NSi{NS-F-P060-016}{q\downarrow0}. The constants may depend on
\NSi{NS-F-P060-017}{X_{\max}}. The formal expansion of
\NSi{NS-F-P060-018}{(u_B,p_B,\mathcal T_{\mathrm{phys}})} has the finite truncations
\NSi{NS-F-P060-019}{\mathcal U^{[N]}} from Proposition~\ref{prop:5.3}, with the physical Stokes
streamfunctions \NSi{NS-F-P060-020}{\mathcal S_n} from \eqref{eq:5.27} summed before differentiation.

Fix radial endpoints \NSi{NS-F-P060-021}{0<X_{\mathrm{lo}}<X_a<X_b<X_{\mathrm{hi}}<\infty}.
The rectangle \NSi{NS-F-P060-022}{[X_{\mathrm{lo}},X_{\mathrm{hi}}]\times[-1,1]} is the fixed radial
enlargement of the closed active annulus
\NSi{NS-F-P060-023}{[X_a,X_b]\times[-1,1]} from Theorem~\ref{thm:4.6}(ii), with
\NSi{NS-F-P060-024}{\eta} ranging over its full closed parameter interval. Uniformly on this enlarged
rectangle, for every fixed composition \NSi{NS-F-P060-025}{\mathscr D^I} of
\NSi{NS-F-P060-026}{q\partial_q,\partial_X,\partial_\eta}, as
\NSi{NS-F-P060-027}{q\downarrow0},
\NSFormula{NS-F-P060-028}{display}{5.42}
\begin{equation}
 \begin{aligned}
  \mathscr D^I(q^Au_{\theta,B}-E_0)&=O_I(q^{2h}),
  &\mathscr D^I(q^Au_{z,B}-U_0)&=O_I(q^{2h}),\\
  \mathscr D^I(q^{1/2}u_{r,B}-V_0/\sqrt{2X})&=O_I(q^{2h}),
  &\mathscr D^I(q^Au_{r,B})&=O_I(q^h).
 \end{aligned}
 \tag{5.42}\label{eq:5.42}
\end{equation}
The constants in these estimates may depend on the fixed radial endpoints. With
\NSi{NS-F-P060-029}{\widehat{\mathcal T}=q^{A+1/2}\mathcal T_{\mathrm{phys}}}, the weighted comparison uses
\NSi{NS-F-P060-030}{\zeta,\delta} from Theorem~\ref{thm:4.6}(iv) and holds on
\NSi{NS-F-P060-031}{X_a<X<X_b}, \NSi{NS-F-P060-032}{-1\leq\eta\leq1}:
\NSFormula{NS-F-P060-033}{display}{5.43}
\begin{equation}
 |\mathscr D^I(\widehat{\mathcal T}-\mathcal T_0)|
 \leq C_Iq^{2h}\zeta\delta^{-N_I}
 \leq C_Iq^h\zeta\delta^{-N_I}
 \qquad(0<q\leq1).
 \tag{5.43}\label{eq:5.43}
\end{equation}
Every fixed physical space--time derivative of
\NSi{NS-F-P060-034}{u_B,p_B,\mathcal T_{\mathrm{phys}},\mathcal E_B} extends continuously to
\NSi{NS-F-P060-035}{\eta=\pm1} on compact sets with \NSi{NS-F-P060-036}{q>0}. The positive-order
corrections \NSi{NS-F-P060-037}{u_B-u^{(0)}} and \NSi{NS-F-P060-038}{p_B-p^{(0)}} vanish for
\NSi{NS-F-P060-039}{X\geq X_+}, with \NSi{NS-F-P060-040}{X_+} from Lemma~\ref{lem:5.2}; in
particular the leading profile's exact exterior \eqref{eq:4.29} is preserved. On the reserved interval
\NSi{NS-F-P060-041}{X\in\mathcal I_{\mathrm{mean}}\subset(X_a,X_v)} from
Theorem~\ref{thm:4.6}(vi), for every \NSi{NS-F-P060-042}{-1\leq\eta\leq1}, the base velocity agrees
exactly with the leading profile:
\NSFormula{NS-F-P060-043}{display}{5.44}
\begin{equation}
 u_B=u^{(0)},
 \qquad u_{z,B}=0,
 \qquad q^Au_{\theta,B}=E_0=c_{\mathrm{patch}}(1+\eta^2)^{-1}X^{-1/2-\lambda}.
 \tag{5.44}\label{eq:5.44}
\end{equation}
Here \NSi{NS-F-P060-044}{u^{(0)}} denotes the velocity determined by the leading profile in
\eqref{eq:4.3}.
\end{proposition}

\begin{proof}
\textbf{Step 1: verify the summation hypotheses and choose the cutoffs.}
We use the potentials \eqref{eq:5.27} and the physical tuples specified before Lemma~\ref{lem:5.4},
keeping order zero intact. The Cartesian formula for \NSi{NS-F-P060-045}{\mathcal A_n} there verifies
regularity at the axis, and Lemma~\ref{lem:5.2} supplies the common supports and all fixed coefficient
derivative bounds. By \eqref{eq:5.26}, since \NSi{NS-F-P060-046}{D<1/2}, the tuple of coefficients in
\eqref{eq:5.27} and \eqref{eq:5.1}, including the stress components, has a reduction in the power of
\NSi{NS-F-P060-047}{q} bounded by \NSi{NS-F-P060-048}{\ell_m=2A+m}. The coefficient constants
absorb the powers of \NSi{NS-F-P060-049}{n} generated by differentiating
\NSi{NS-F-P060-050}{q^{\lambda_n}}. For the cutoffs, we set
\NSi{NS-F-P060-051}{\sigma=c_nq}. For every integer \NSi{NS-F-P060-052}{j\geq0},
\NSFormula{NS-F-P060-053}{display}{none}
\[
 (q\partial_q)^j\chi(c_nq)
 =(\sigma\partial_\sigma)^j\chi(\sigma)\big|_{\sigma=c_nq},
\]
\NSFormula{NS-F-P060-054}{display}{none}
\[
 \sup_{n\geq1,\,q>0}|(q\partial_q)^j\chi(c_nq)|
 \leq\sup_{\sigma>0}|(\sigma\partial_\sigma)^j\chi(\sigma)|
 =:C_j<\infty.
\]
